\chapter{Lower back pain}

\section*{Definition and Overview}
Lower back pain (LBP) refers to pain, muscle tension, or stiffness localised between the lower margin of the 12th ribs and the inferior gluteal folds, with or without referred leg pain (sciatica). It is a symptom rather than a specific disease entity, arising from a complex interplay of anatomical structures, including the lumbar vertebrae, intervertebral discs, facet joints, spinal cord and nerves, and surrounding paraspinal muscles and ligaments. LBP is clinically classified by duration as acute (lasting less than 6 weeks), subacute (6 to 12 weeks), or chronic (12 weeks or more). In over 90\% of cases, it is classified as non-specific lower back pain, meaning no clear anatomical origin or serious underlying pathology can be identified. Although LBP is generally self-limiting, it represents a leading cause of functional disability, absenteeism, and healthcare utilisation worldwide.

\section*{Epidemiology}
Lower back pain is a ubiquitous health condition with an extraordinary global burden. It is estimated that approximately 80\% of individuals will experience at least one episode of LBP during their lifetime. According to the Australian Institute of Health and Welfare (AIHW), back pain is one of the most common reasons for visiting a general practitioner and is the leading cause of occupational disability in Australia. The point prevalence of chronic LBP is estimated at 7.3\% globally, peaking in the 40--60 age group. It affects both sexes, though women report slightly higher rates of chronic back pain than men. Key risk factors associated with LBP include physical inactivity, obesity, poor ergonomics, repetitive lifting of heavy loads, psychological factors (such as anxiety, depression, and poor coping strategies), and tobacco smoking, which is linked to accelerated intervertebral disc degeneration.

\section*{Aetiology and Pathophysiology}
The pathomechanisms of lower back pain involve mechanical, inflammatory, neuropathic, and psychosocial factors.
\begin{itemize}
    \item \textbf{Muscle or ligament strain:} Acute micro-tears in the paraspinal muscles or erector spinae ligaments due to sudden movement, lifting trauma, or cumulative postural fatigue.
    \item \textbf{Intervertebral disc herniation:} Protrusion or extrusion of the nucleus pulposus through a tear in the annulus fibrosus, causing mechanical compression and chemical irritation of the adjacent nerve roots (radiculopathy).
    \item \textbf{Facet joint arthropathy:} Osteoarthritic degeneration of the lumbar facet joints, leading to localised mechanical pain and stiffness.
    \item \textbf{Lumbar spinal stenosis:} Narrowing of the central spinal canal or lateral recesses, resulting in compression of the cauda equina nerve roots, typically presenting as neurogenic claudication.
    \item \textbf{Spondylolisthesis:} The forward slippage of one lumbar vertebra over another, causing mechanical instability and potential nerve root impingement.
    \item \textbf{Spondyloarthropathies:} Inflammatory autoimmune conditions (e.g. ankylosing spondylitis) characterised by chronic sacroiliitis and spinal enthesitis.
\end{itemize}

\section*{Clinical Features}
The clinical presentation of lower back pain depends on its underlying mechanical, neuropathic, or inflammatory drivers.
\begin{itemize}
    \item \textbf{Localised lumbar pain:} A dull, aching, or sharp discomfort in the lower back, often exacerbated by movement, bending, or prolonged sitting.
    \item \textbf{Radicular pain (sciatica):} Sharp, shooting, or electric shock-like pain radiating down the posterior or lateral aspect of the leg, often extending below the knee in a dermatomal pattern.
    \item \textbf{Stiffness and restricted mobility:} Reduced range of motion in the lumbar spine, particularly during flexion, extension, or lateral rotation.
    \item \textbf{Muscle spasms:} Involuntary, painful contractions of the paraspinal muscles, causing postural deviation or guarding.
    \item \textbf{Sensory and motor deficits:} Numbness, tingling, or weakness in specific dermatomal or myotomal distributions of the lower limb (e.g. foot drop).
\end{itemize}

\section*{Differential Diagnosis}
It is critical to distinguish benign mechanical lower back pain from serious spinal pathologies (referred to as "red flags") and non-spinal sources of referred pain.
\begin{itemize}
    \item \textbf{Cauda equina syndrome:} Compression of the lumbosacral nerve roots, presenting as saddle anaesthesia, bladder or bowel dysfunction, and progressive bilateral motor weakness.
    \item \textbf{Spinal malignancy:} Primary bone tumours or metastatic lesions (e.g. from breast, prostate, or lung cancer) presenting as progressive night pain unresponsive to rest.
    \item \textbf{Vertebral osteomyelitis or discitis:} Infection of the spine or intervertebral discs, characterised by localised spinal tenderness, fever, and elevated inflammatory markers.
    \item \textbf{Vertebral compression fracture:} Skeletal failure due to osteoporosis or trauma, presenting as sudden onset of severe, localised focal pain.
    \item \textbf{Abdominal aortic aneurysm (AAA):} A vascular emergency presenting as deep, boring back pain, potentially accompanied by a pulsatile abdominal mass.
    \item \textbf{Nephrolithiasis (kidney stones) or pyelonephritis:} Renal pathologies presenting as acute loin-to-groin pain and costovertebral angle tenderness.
\end{itemize}

\section*{When to Seek Medical Care}
Most acute episodes of lower back pain improve with conservative care. However, patients should seek medical advice if the pain is severe, persistent beyond a few weeks, or associated with constitutional symptoms. Immediate emergency evaluation is mandatory if red flag symptoms indicating cauda equina syndrome or other life-threatening conditions are present.

\textbf{Call 000 (or local emergency services) for:}
\begin{itemize}
    \item \textbf{Saddle anaesthesia:} Loss of sensation in the perineum, buttocks, inner thighs, or perianal region.
    \item \textbf{Bladder or bowel dysfunction:} New onset of urinary retention, overflow incontinence, or faecal incontinence.
    \item \textbf{Progressive motor weakness:} Sudden or worsening difficulty lifting the foot (foot drop), walking, or moving the legs.
    \item \textbf{Sudden severe back and abdominal pain:} Tearing, severe pain that radiates to the groin, which may indicate a ruptured abdominal aortic aneurysm.
\end{itemize}

\section*{Diagnosis and Investigations}
Routine imaging is not recommended for non-specific lower back pain, as incidental findings (such as disc bulges) are highly prevalent in asymptomatic individuals. Diagnostic tests are reserved for patients with suspected serious pathology or persistent radicular symptoms.
\begin{itemize}
    \item \textbf{Clinical evaluation:} Detailed history and physical examination, including straight leg raise (SLR) test, neurological assessment of reflexes, sensation, and motor power.
    \item \textbf{Magnetic resonance imaging (MRI):} The imaging modality of choice for suspected disc herniation, spinal stenosis, infection, or malignancy.
    \item \textbf{Lumbar spine X-ray:} Useful for evaluating structural bony abnormalities, fractures, or spondylolisthesis, but has limited utility for soft tissue evaluation.
    \item \textbf{CT scan:} Recommended when MRI is contraindicated or to obtain detailed bony architecture in complex fractures.
    \item \textbf{Inflammatory markers (ESR, CRP):} Indicated if spinal infection or systemic inflammatory arthritis (e.g. spondyloarthropathy) is suspected.
\end{itemize}

\section*{Management}
The management of lower back pain focuses on reassuring the patient, keeping them active, and avoiding prolonged bed rest.
\begin{itemize}
    \item \textbf{Encouraging activity:} Recommending that patients remain active and continue daily activities as tolerated, rather than resting in bed.
    \item \textbf{Pharmacological management:} First-line analgesics such as paracetamol or oral non-steroidal anti-inflammatory drugs (NSAIDs) for short-term pain relief; muscle relaxants may be used briefly for severe spasms.
    \item \textbf{Physiotherapy:} Structured exercise programmes focusing on core stabilisation, stretching, general aerobic exercise, and posture correction.
    \item \textbf{Psychosocial support (CBT):} Cognitive behavioural therapy to address negative coping patterns, fear-avoidance behaviours, and chronic pain management.
    \item \textbf{Interventional therapies:} Epidural steroid injections or facet joint blocks for selected patients with severe, persistent radiculopathy or arthropathy.
    \item \textbf{Surgical consultation:} Reserved for patients with progressive neurological deficits, cauda equina syndrome, or intractable pain unresponsive to non-operative treatment.
\end{itemize}

\section*{Complications}
\begin{itemize}
    \item \textbf{Transition to chronic pain:} Approximately 10\% of acute LBP cases develop into chronic, disabling pain with long-term psychological impacts.
    \item \textbf{Permanent neurological deficit:} Chronic nerve root compression leading to persistent muscle weakness, foot drop, or sensory loss.
    \item \textbf{Functional disability and loss of employment:} Significant impairment in physical function, limiting work capacity and daily activities.
    \item \textbf{Medication side effects:} Gastrointestinal bleeding, renal dysfunction, or cardiovascular risks from chronic NSAID use, or dependence on opioid analgesics.
    \item \textbf{Mental health issues:} Strong correlation between chronic back pain and the development of clinical depression and anxiety.
\end{itemize}

\section*{Prognosis and Outcomes}
The prognosis for acute lower back pain is highly favourable. The majority of patients (60\% to 70\%) recover substantially within 6 weeks, and up to 90\% recover within 12 weeks. Recurrence is common, however, with approximately half of patients experiencing another episode within a year. The prognosis is less favourable for individuals with high levels of baseline pain, radicular symptoms, or significant psychosocial distress (often referred to as "yellow flags"). For chronic LBP, outcomes are best when managed with a multidisciplinary, active rehabilitation approach focused on functional improvement and self-management, rather than passive medical interventions.

\section*{Prevention and Self-Care}
\begin{itemize}
    \item \textbf{Maintain regular physical activity:} Engaging in low-impact aerobic exercise, such as walking or swimming, and core strengthening.
    \item \textbf{Practice ergonomic lifting:} Bending the knees and keeping the back straight when lifting objects, and avoiding twisting under load.
    \item \textbf{Optimize workstation setup:} Ensuring proper chair height, lumbar support, and screen alignment to maintain neutral spinal alignment.
    \item \textbf{Maintain a healthy weight:} Reducing mechanical strain on the lumbar spine by managing body weight.
    \item \textbf{Avoid prolonged sitting:} Taking regular breaks to stand, stretch, and walk during prolonged desk work.
\end{itemize}

\section*{Sources and Further Reading}
The following reputable organisations provide further information and clinical guidelines on lower back pain management:
\begin{itemize}
    \item Healthdirect Australia. \textit{Lower back pain symptoms, causes, and treatment options.} https://www.healthdirect.gov.au/
    \item Mayo Clinic. \textit{Comprehensive overview of back pain diagnosis and treatment.} https://www.mayoclinic.org/
    \item National Institute for Health and Care Excellence. \textit{LBP and sciatica in over 16s: assessment and management.} https://www.nice.org.uk/
    \item Arthritis Australia. \textit{Consumer resources on managing back pain and joint health.} https://arthritisaustralia.com.au/
    \item World Health Organization. \textit{Global guidelines on chronic primary low back pain.} https://www.who.int/
\end{itemize}
